\section{Ergebnis und Ausblick}

Das Modell zeigt, dass Kreis und Ellipse nicht an Sinus und Kosinus als
einzige Parametrisierung gebunden sind. Ihre eigentliche Invariante ist eine
algebraische Beziehung zwischen zwei Koordinaten. Die Dreieckswelle und der
Kosinus erscheinen dadurch als zwei Grenzfälle derselben Familie von
Zeitordnungen.

Gate~2 schließt die Summierbarkeit, Nennerpositivität und globale
Beschränktheit der geglätteten Steuerung. Die nächsten formalen Schritte sind:
\begin{enumerate}
  \item Identifikation der Randfälle mit Dreieckswelle und Kosinus,
  \item Glattheit für $0<\lambda$ und termweise Differenzierbarkeit,
  \item Berechnung der endlichen Umkehrgeschwindigkeit für $\lambda>0$,
  \item Verallgemeinerung auf höhere Sphären und Ellipsoide.
\end{enumerate}

Besonders interessant ist dabei die letzte Perspektive: Der zweidimensionale
Wurzellift ist möglicherweise nur der einfachste Fall eines allgemeineren
Prinzips, nach dem verschiedene Dynamiken dieselbe quadratische geometrische
Invariante realisieren.
